\section{The growing-band position profile}
\label{sec:terminal-modal-position-profile}

The first scalar target in
\eqref{eq:terminal-modal-source-position-target} can be closed without any
velocity estimate.  The proof keeps the five exact pieces from the end of
\cref{sec:terminal-modal-admission} separate.  Put
\begin{equation}
 H_m=\left\lfloor\sqrt{\frac{m}{64\log(16m)}}\right\rfloor,
 \qquad \phi_s=\frac{s\pi}{m},
 \label{eq:terminal-modal-bandwidth}
\end{equation}
and write
\begin{equation}
 Y_{m+1}(2\phi_s)-mG_{2s}
 =B_{m,s}+R^0_{m,s}+R^\Delta_{m,s}+R^M_{m,s}+E_{m,s}.
 \label{eq:terminal-modal-position-five-pieces}
\end{equation}
Here $B_{m,s}$ is the homogeneous and $n<4$ response minus the ideal packet,
$R^0_{m,s}$ is the response to the constant derivative trace
$U_0=3q^3/40$, $R^\Delta_{m,s}$ is the response to $U_n-U_0$,
$R^M_{m,s}$ is the response to $M_n$, and $E_{m,s}$ is the final
degree/endpoint correction.  This is the exact linear split already defined
after \cref{lem:terminal-modal-source-scalar-reduction}; no limiting model is
being substituted in \eqref{eq:terminal-modal-position-five-pieces}.

Define, for $0\le x\le1$,
\begin{equation}
 p_0(x)=2\frac{2t_x+t_x^2/5-1/5}{1+t_x^2},
 \qquad
 t_x=e^{-x/8},
 \qquad
 u(x)=\frac{p_0(x)-4/5}{16},
 \label{eq:terminal-modal-position-limit-profile}
\end{equation}
so $u(0)=3/40$.  The following lemma records all uniform remainder estimates
needed below.

\begin{lemma}[Uniform position source limits; proved here]
\label{lem:terminal-modal-position-source-limits}
Uniformly for $1\le s\le H_m$, as $m\to\infty$,
\begin{align}
 \frac{B_{m,s}}{q^2}&=o(1),
 &\frac{E_{m,s}}{q^2}&=o(1),
 \label{eq:terminal-modal-position-small-pieces}\\
 \frac{R^0_{m,s}}{q^2}
 &=I^0_s+o(1),
 &
 I^0_s&=(-1)^{s+1}\int_0^1 e^{-(1-x)/16}
 \frac3{40}\sin(2\pi sx)\sin(\pi sx)\,dx,
 \label{eq:terminal-modal-position-constant-limit}\\
 \frac{R^\Delta_{m,s}}{q^2}
 &=I^\Delta_s+o(1),
 &
 I^\Delta_s&=(-1)^{s+1}\int_0^1 e^{-(1-x)/16}
 \left(u(x)-\frac3{40}\right)
 \sin(2\pi sx)\sin(\pi sx)\,dx,
 \label{eq:terminal-modal-position-varying-limit}
\end{align}
and
\begin{equation}
 \frac{|R^M_{m,s}|}{q^2}
 \le \frac{43}{10240\pi s}+o(1).
 \label{eq:terminal-modal-position-mass-limit}
\end{equation}
Every $o(1)$ in these displays is uniform over the growing band.
\end{lemma}

\begin{proof}
We first derive the common derivative-trace limit.  If $x_n=n/m$, then
\begin{equation}
 \chi^{n-1}=e^{-x_n/8}\bigl(1+O(q)\bigr)
 \label{eq:terminal-modal-position-chi-limit}
\end{equation}
uniformly for $4\le n\le m-1$.  Indeed
$\log\chi=-2q+O(q^3)$ and $mq=1/16$.  Substitution in
\eqref{eq:terminal-modal-frontier-optimum} and then
\eqref{eq:terminal-modal-source-U} gives
\begin{equation}
 \frac{U_n}{q^3}=u(x_n)+O(q)
 \label{eq:terminal-modal-position-U-limit}
\end{equation}
uniformly on the same range.

For $z=e^{2i\phi_s}$, direct Taylor expansion of the exact source factor is
\begin{equation}
 2\Re\!\left[z^{n-2}(1-z)(1+3z)\right]
 =16\phi_s\sin(2\pi s x_n)+O(\phi_s^2),
 \label{eq:terminal-modal-position-source-Taylor}
\end{equation}
uniformly in $n$.  The Green factor in
\eqref{eq:terminal-modal-source-green} satisfies
\begin{align}
 r^{m-n}&=e^{-(1-x_n)/16}
 \left(1+O(q+m\phi_s^2)\right),
 \label{eq:terminal-modal-position-damping-Taylor}\\
 \sin((m-n+1)\phi_s)
 &=(-1)^{s+1}\sin((n-1)\phi_s),
 \qquad
 \frac{\phi_s}{\sin\phi_s}=1+O(\phi_s^2).
 \label{eq:terminal-modal-position-Green-Taylor}
\end{align}
On the band, $m\phi_s^2=O(1/\log(16m))=o(1)$.  Combining
\eqref{eq:terminal-modal-position-U-limit}--
\eqref{eq:terminal-modal-position-Green-Taylor}, using
$16mq=1$, and replacing the resulting Riemann sum by its integral gives
\begin{equation}
 \frac{R^U_{m,s}}{q^2}
 =(-1)^{s+1}\int_0^1e^{-(1-x)/16}u(x)
 \sin(2\pi sx)\sin(\pi sx)\,dx
 +O\!\left(q+\frac{s}{m}+\frac{s^2}{m}\right).
 \label{eq:terminal-modal-position-U-Riemann}
\end{equation}
The Riemann error is uniform because the smooth amplitude has bounded
derivative and the trigonometric factor has derivative $O(s)$.  The same
calculation with $u(x)$ replaced by the constant $3/40$ proves
\eqref{eq:terminal-modal-position-constant-limit}; subtraction proves
\eqref{eq:terminal-modal-position-varying-limit}.

For completeness, the constant response also has an exact finite form.  If
$c=\cos\phi_s$, $x=(1-q)cz^{-1}$, and $K=m-4$, then
\begin{align}
 R^0_{m,s}
 &=2U_0\Re\!\left[z^{-2}(1-z)(1+3z)T_K\right],
 \label{eq:terminal-modal-position-constant-exact}\\
 T_K
 &=\frac{2cx-x^2-U_{K+1}(c)x^{K+1}+U_K(c)x^{K+2}}
 {1-2cx+x^2},
 \label{eq:terminal-modal-position-Chebyshev-sum}
\end{align}
where $U_k(c)=\sin((k+1)\phi_s)/\sin\phi_s$.  This is the finite geometric
Chebyshev identity $T_K=\sum_{\ell=1}^KU_\ell(c)x^\ell$; at the integer mode,
\begin{equation}
 U_K(c)=(-1)^{s+1}\frac{\sin(3\phi_s)}{\sin\phi_s},
 \qquad
 U_{K+1}(c)=(-1)^{s+1}\frac{\sin(2\phi_s)}{\sin\phi_s}.
 \label{eq:terminal-modal-position-Chebyshev-end}
\end{equation}
Thus the finite cutoff and the reflected endpoint are already included in
\eqref{eq:terminal-modal-position-constant-limit}.

The sharper mass estimate
\eqref{eq:terminal-modal-finite-mass-bound}, together with the identification
$\mathfrak m_n=M_n/q^4$ after the proved chronology, gives
$0<M_n<43q^4/80$.  Taking absolute values only for this smaller term in
\eqref{eq:terminal-modal-source-green} yields
\begin{equation}
 \frac{|R^M_{m,s}|}{q^2}
 \le\frac{43}{40}\frac{q^2m}{\sin\phi_s}
 =\frac{43}{10240\,m\sin\phi_s}.
 \label{eq:terminal-modal-position-mass-finite}
\end{equation}
Since $m\sin\phi_s=\pi s(1+O(\phi_s^2))$ uniformly on the band, this proves
\eqref{eq:terminal-modal-position-mass-limit}.

It remains to justify the two small pieces.  Let $\mathcal A_n$ denote the
even transform from \eqref{eq:terminal-modal-driven-transform}.  Propagation
from the two fixed early values gives exactly
\begin{equation}
 B^{\rm raw}_{m,s}
 =\frac{r^{m-3}}{\sin\phi_s}
 \left\{\mathcal A_4\sin((m-2)\phi_s)
 -r\mathcal A_3\sin((m-3)\phi_s)\right\}.
 \label{eq:terminal-modal-position-base-exact}
\end{equation}
The reflected ideal packet and the fixed-prefix correction recurrence give,
for $n=3,4$,
\begin{equation}
 \mathcal A_n=-q^2(1-q)^{n-1}\cos^{n-1}\phi_s
 \cos((n-1)\phi_s)+O(q^3).
 \label{eq:terminal-modal-position-base-data}
\end{equation}
The $O(q^3)$ is uniform because only finitely many correction coordinates are
involved.  There is no small-sine amplification of this error:
\begin{equation}
 \left|\frac{r^{m-3}}{\sin\phi_s}
 \left\{O(q^3)\sin(2\phi_s)-rO(q^3)\sin(3\phi_s)\right\}\right|
 \le 5q^3O(1),
 \label{eq:terminal-modal-position-base-error}
\end{equation}
because $|\sin(k\phi)/\sin\phi|\le k$.  Substitution of the ideal terms in
\eqref{eq:terminal-modal-position-base-exact}, followed by
$\cos(2\phi)\sin(3\phi)-\cos(3\phi)\sin(2\phi)=\sin\phi$, gives
\begin{equation}
 B^{\rm raw}_{m,s}=(-1)^{s+1}q^2r^m+O(q^3).
 \label{eq:terminal-modal-position-base-match}
\end{equation}
The ideal packet transform is explicitly
\begin{equation}
 mG_{2s}=(-1)^{s+1}q^2r^m+q^2(1-q)^m2^{-m}.
 \label{eq:terminal-modal-position-ideal-transform}
\end{equation}
Thus the main terms cancel, and the exponentially small endpoint atom proves
the first estimate in \eqref{eq:terminal-modal-position-small-pieces}.

Finally let $\sigma_n$ be the raw frontier value after the proper-prefix
step.  The exact frontier recurrence used in
\cref{lem:terminal-modal-chronology-reduction}, together with
\eqref{eq:terminal-modal-frontier-bounds}, gives
\begin{equation}
 |\sigma_n|\le\frac12|\sigma_{n-1}|+O(q^3),
 \qquad
 |\sigma_m|+|\sigma_{m-1}|=O(q^3)+O(2^{-m}q^2).
 \label{eq:terminal-modal-position-frontier-decay}
\end{equation}
Using $P_m=\Delta_mf_{m-1}$,
$P_{m-1}=\Delta_mf_m+2q^3/(5\eta)$, and
$f_m+f_{m-2}=2af_{m-1}/\eta$, the final endpoint entry
\eqref{eq:terminal-modal-source-final-endpoint} simplifies exactly to
\begin{equation}
 W_m^{\rm end}
 =\frac{1-q}{2}A_m^{\rm fr}
 -\frac{q(1-q)}2P_m
 +\frac{q\eta(1-q)}4P_{m-1}
 +\frac{q^3(1+q)}{10}.
 \label{eq:terminal-modal-position-endpoint-cancellation}
\end{equation}
Here $A_m^{\rm fr}=-\eta\sigma_{m-1}/2+q^3/5$, while the literal new endpoint
residual subtracted in $Y_{m+1}$ is $-\eta\sigma_m+q^3/5$.  Therefore both
are $O(q^3)$ by \eqref{eq:terminal-modal-position-frontier-decay}; the same is
true of \eqref{eq:terminal-modal-position-endpoint-cancellation} and of the
two remaining final entries $U_m,2U_m$.  Hence $E_{m,s}=O(q^3)$ uniformly in
$s$, proving the second estimate in
\eqref{eq:terminal-modal-position-small-pieces}.
\end{proof}

\begin{theorem}[Growing-band entry position profile; proved here]
\label{lem:terminal-modal-position-profile}
For all sufficiently large $m$, the literal full-face entry coefficients
satisfy
\begin{equation}
 \frac m\alpha|C_{2s}-G_{2s}|\le\frac1{256},
 \qquad 1\le s\le H_m.
 \label{eq:terminal-modal-entry-position-profile}
\end{equation}
\end{theorem}

\begin{proof}
Put $a_0=1/16$.  Product-to-sum and direct integration give
\begin{equation}
 I^0_s=\frac3{80}(-1)^{s+1}(J_s-J_{3s}),
 \qquad
 J_k=\frac{a_0((-1)^k-e^{-a_0})}{a_0^2+k^2\pi^2}.
 \label{eq:terminal-modal-position-constant-integral}
\end{equation}
If $s$ is odd, use $1+e^{-a_0}<2$ and take the difference of the two
denominators; if $s$ is even, use $1-e^{-a_0}<a_0$.  In both cases,
\begin{equation}
 |I^0_s|\le\frac1{240\pi^2s^2}.
 \label{eq:terminal-modal-position-constant-bound}
\end{equation}

The function $p_0(x)$ in
\eqref{eq:terminal-modal-position-limit-profile} is decreasing in $x$, so
$u(x)$ is decreasing.  Hence
\begin{equation}
 f(x)=e^{-(1-x)/16}\left(u(x)-\frac3{40}\right)
 \label{eq:terminal-modal-position-varying-amplitude}
\end{equation}
decreases from zero and has total variation $3/40-u(1)$.  The alternating
Taylor bound gives
\begin{equation}
 e^{-1/8}>\frac{2711}{3072},
 \qquad
 2-p_0(1)<2-p_0\!\left(-8\log\frac{2711}{3072}\right)
 =\frac{5478536}{83933525}<\frac1{15}.
 \label{eq:terminal-modal-position-variation-rational}
\end{equation}
Equivalently, in the middle expression we evaluate the increasing rational
function $2(2t+t^2/5-1/5)/(1+t^2)$ at
$t=2711/3072$.  Thus $\operatorname{TV}(f)<1/240$.
Integration by parts after
$\sin(2y)\sin y=(\cos y-\cos3y)/2$ now gives
\begin{equation}
 |I^\Delta_s|
 \le\frac{2\operatorname{TV}(f)}{3\pi s}
 <\frac1{360\pi s}.
 \label{eq:terminal-modal-position-varying-bound}
\end{equation}

By \cref{lem:terminal-modal-position-source-limits} and $s\ge1$,
\begin{equation}
 \frac{|Y_{m+1}(2\phi_s)-mG_{2s}|}{q^2}
 \le \frac1{240\pi^2}+\frac1{360\pi}
 +\frac{43}{10240\pi}+o(1)
 <\frac{257}{92160}+o(1)<\frac1{256}
 \label{eq:terminal-modal-position-final-ledger}
\end{equation}
uniformly on the band.  The penultimate inequality uses $\pi>3$; the final
strict rational slack is
$1/256-257/92160=103/92160$.  Since $\alpha=q^2$ and
$mC_{2s}=Y_{m+1}(2\phi_s)$, this is
\eqref{eq:terminal-modal-entry-position-profile} for all sufficiently large
$m$.
\end{proof}
